
% ============================================================
\subsection{LoMFJ}
% ============================================================

LoMFJ fungiert als Schnittstelle zwischen der Hauptapplikation und externen Kommunikationssystemen, beispielsweise Telefonie- oder Alarmierungssystemen.  
Dabei erzeugt LoMFJ eigenständig Logdateien über stattfindende Ereignisse und Fehler, unabhängig vom zentralen Logging der Hauptapplikation.

\subsubsection{Setup}

Im Gegensatz zur Hauptapplikation, deren Logging zentralisiert ist, schreibt jede LoMFJ-Node eigene lokale Logdateien.

Für die Konfiguration wird eine einzelne Datei verwendet, die sowohl die \textit{RollingFile}-Appender als auch die \textit{Logger}-Definitionen enthält:

\begin{quote}
    Konfiguration des Loggings: \textit{~/server/lomfj/config/log4j2.xml}
\end{quote}

Das Logging ist integraler Bestandteil des LoMFJ-Services und kann nicht unabhängig gestartet oder gestoppt werden.

\subsubsection{Logs}

Abhängig von der kundenspezifischen MFJ-Version können unterschiedliche Logdateien vorhanden sein.  
Im Folgenden sind die Logbereiche aufgeführt, wie sie beim Kunden \texttt{SNDE} eingesetzt werden.

\paragraph{\lomfjAuditLog}
Enthält Informationen über Benutzeranmeldungen.

\paragraph{\lomfjDfLog}
Protokolliert Aktionen, die durch Interaktionen zwischen MFJ und den DF-Steckern ausgelöst werden.

\paragraph{\lomfjErrorLog}
Sammelt Fehlerzustände und Ausnahmen, die bei Interaktionen zwischen MFJ und angebundenen Komponenten auftreten.

\paragraph{\lomfjEventmanagerLog}
Der Event Manager verarbeitet eingehende Informationen aus angebundenen Services (z.\,B. FNAS).  
Er erkennt beispielsweise doppelte Call-PDUs und stellt sicher, dass auf Pelix-Seite nur ein konsolidiertes \texttt{CallEvent} erzeugt wird.

\paragraph{\lomfjFnasLog}
Protokolliert sämtliche Interaktionen mit dem FNAS.

\paragraph{\lomfjHeartbeatLog}
Damit primäre und sekundäre MFJ-Node zusammenarbeiten können, wird regelmäßig ein Heartbeat erzeugt, der den aktiven Zustand der Node signalisiert. Diese Informationen werden in diesem Log festgehalten.

\paragraph{\lomfjIncompleteFnasPdusLog}
Enthält Requests vom FNAS, die unvollständige oder fehlerhafte Felder enthalten.

\paragraph{\lomfjMfjLog}
Beinhaltet allgemeine Systeminformationen zur MFJ-Architektur, wie Speicherverbrauch, Startzeitpunkt oder Prozess-ID.  
Diese Logs sind im Normalbetrieb meist wenig relevant, können jedoch bei ungewöhnlichem Verhalten wichtige Hinweise liefern (z.\,B. SDS-Protokolle).

\paragraph{\lomfjMfjServerLog}
Für den Kunden \texttt{SNDE} wurde gefordert, alle LoMFJ-Logeinträge zusätzlich in einer zentralen Datei zu erfassen.  
Diese Datei entspricht funktional einem \texttt{server.log} und enthält sämtliche Ereignisse in komprimierter Form.

Aufgrund des hohen Logaufkommens wachsen diese Dateien sehr schnell an und müssen bei Bedarf zeitnah ausgewertet werden.  
Diese Logs enthalten unter anderem Status- und Ereignismeldungen von Sensoren und fest installierten Alarmkomponenten.

\paragraph{\lomfjPelixLog}
Protokolliert die Interaktionen zwischen MFJ und Pelix bzw. mLogistics.

\paragraph{\lomfjUeaLog}
Betrifft die ÜEA-Schnittstelle (Überfall- und Einbruchmeldeanlagen) gemäß den VdS-Richtlinien.

\paragraph{\EMPerformance}
Betrifft die ÜEA-Schnittstelle (Überfall- und Einbruchmeldeanlagen) gemäß den VdS-Richtlinien.

\paragraph{\PelixPerformance}
Betrifft die ÜEA-Schnittstelle (Überfall- und Einbruchmeldeanlagen) gemäß den VdS-Richtlinien.
